% Shared preamble for the course decision diagrams.
% Each diagram file starts with % Shared preamble for the course decision diagrams.
% Each diagram file starts with % Shared preamble for the course decision diagrams.
% Each diagram file starts with \input{diagram_style} and contains one tikzpicture.
% Compile with pdflatex; build_diagrams.py then crops each page to its content.
\documentclass[11pt]{article}
\usepackage[a3paper,landscape,margin=10mm]{geometry}
\usepackage[T1]{fontenc}
\usepackage{tgheros}                 % Neudata body font (as in neudata.cls)
\renewcommand{\sfdefault}{qhv}
\renewcommand{\familydefault}{\sfdefault}
\usepackage[scaled=0.95,varqu]{zi4}
\usepackage{tikz}
\usetikzlibrary{arrows.meta,positioning,fit,calc}
\pagestyle{empty}
\setlength{\parindent}{0pt}
\hyphenpenalty=10000
\exhyphenpenalty=10000

% Neudata palette (neu-data/.github BRAND.md); names kept so the diagram sources need no edits
\definecolor{cblue}{HTML}{0B376C}    % Neudata Blue: question boxes, links between diagrams
\definecolor{corange}{HTML}{B5652B}  % muted contrasting colour: warnings
\definecolor{cgreen}{HTML}{055F56}   % Neudata Teal: method boxes
\definecolor{clight}{HTML}{EAF2F5}   % Neudata Mist: question box fill
\definecolor{ctext}{HTML}{1F2D33}    % Neudata Ink: text
\definecolor{cgrey}{HTML}{8A99A3}    % Neudata Grey: further-learning routes
\definecolor{cpale}{HTML}{F4F8FA}
\definecolor{cnavy}{HTML}{04242F}    % Deep Navy: titles

\tikzset{
  every node/.style={text=ctext, font=\small},
  title/.style={font=\Large\bfseries, text=cnavy, anchor=west},
  subtitle/.style={font=\normalsize, text=ctext!72, anchor=west},
  q/.style={draw=cblue, fill=clight, rounded corners=3pt, line width=0.8pt,
            align=center, inner sep=6pt, text width=#1},
  q/.default=5cm,
  m/.style={draw=cgreen, fill=white, rounded corners=3pt, line width=1pt,
            align=left, inner sep=6pt, text width=#1},
  m/.default=5cm,
  opt/.style={draw=cblue, fill=white, rounded corners=3pt, line width=0.6pt,
              align=center, inner sep=5pt, text width=#1},
  opt/.default=3.2cm,
  further/.style={draw=cgrey, dashed, fill=white, rounded corners=3pt, line width=0.8pt,
                  align=left, inner sep=6pt, text=ctext, text width=#1},
  further/.default=5cm,
  warn/.style={draw=corange, fill=white, rounded corners=3pt, line width=1pt,
               align=left, inner sep=6pt, text width=#1},
  warn/.default=6cm,
  note/.style={align=left, text width=#1, font=\footnotesize, text=ctext},
  note/.default=5cm,
  band/.style={font=\normalsize\bfseries, text=cgreen, anchor=west, align=left},
  flow/.style={-{Stealth[length=2.6mm]}, line width=0.8pt, draw=ctext!75},
  flowg/.style={-{Stealth[length=2.6mm]}, line width=0.8pt, draw=cgrey, dashed},
  lab/.style={font=\footnotesize, fill=white, inner sep=2pt, text=ctext},
  code/.style={font=\ttfamily\footnotesize},
}
\newcommand{\rc}[1]{{\ttfamily\footnotesize #1}}
\newcommand{\dref}[1]{{\color{cblue}#1}}
 and contains one tikzpicture.
% Compile with pdflatex; build_diagrams.py then crops each page to its content.
\documentclass[11pt]{article}
\usepackage[a3paper,landscape,margin=10mm]{geometry}
\usepackage[T1]{fontenc}
\usepackage{tgheros}                 % Neudata body font (as in neudata.cls)
\renewcommand{\sfdefault}{qhv}
\renewcommand{\familydefault}{\sfdefault}
\usepackage[scaled=0.95,varqu]{zi4}
\usepackage{tikz}
\usetikzlibrary{arrows.meta,positioning,fit,calc}
\pagestyle{empty}
\setlength{\parindent}{0pt}
\hyphenpenalty=10000
\exhyphenpenalty=10000

% Neudata palette (neu-data/.github BRAND.md); names kept so the diagram sources need no edits
\definecolor{cblue}{HTML}{0B376C}    % Neudata Blue: question boxes, links between diagrams
\definecolor{corange}{HTML}{B5652B}  % muted contrasting colour: warnings
\definecolor{cgreen}{HTML}{055F56}   % Neudata Teal: method boxes
\definecolor{clight}{HTML}{EAF2F5}   % Neudata Mist: question box fill
\definecolor{ctext}{HTML}{1F2D33}    % Neudata Ink: text
\definecolor{cgrey}{HTML}{8A99A3}    % Neudata Grey: further-learning routes
\definecolor{cpale}{HTML}{F4F8FA}
\definecolor{cnavy}{HTML}{04242F}    % Deep Navy: titles

\tikzset{
  every node/.style={text=ctext, font=\small},
  title/.style={font=\Large\bfseries, text=cnavy, anchor=west},
  subtitle/.style={font=\normalsize, text=ctext!72, anchor=west},
  q/.style={draw=cblue, fill=clight, rounded corners=3pt, line width=0.8pt,
            align=center, inner sep=6pt, text width=#1},
  q/.default=5cm,
  m/.style={draw=cgreen, fill=white, rounded corners=3pt, line width=1pt,
            align=left, inner sep=6pt, text width=#1},
  m/.default=5cm,
  opt/.style={draw=cblue, fill=white, rounded corners=3pt, line width=0.6pt,
              align=center, inner sep=5pt, text width=#1},
  opt/.default=3.2cm,
  further/.style={draw=cgrey, dashed, fill=white, rounded corners=3pt, line width=0.8pt,
                  align=left, inner sep=6pt, text=ctext, text width=#1},
  further/.default=5cm,
  warn/.style={draw=corange, fill=white, rounded corners=3pt, line width=1pt,
               align=left, inner sep=6pt, text width=#1},
  warn/.default=6cm,
  note/.style={align=left, text width=#1, font=\footnotesize, text=ctext},
  note/.default=5cm,
  band/.style={font=\normalsize\bfseries, text=cgreen, anchor=west, align=left},
  flow/.style={-{Stealth[length=2.6mm]}, line width=0.8pt, draw=ctext!75},
  flowg/.style={-{Stealth[length=2.6mm]}, line width=0.8pt, draw=cgrey, dashed},
  lab/.style={font=\footnotesize, fill=white, inner sep=2pt, text=ctext},
  code/.style={font=\ttfamily\footnotesize},
}
\newcommand{\rc}[1]{{\ttfamily\footnotesize #1}}
\newcommand{\dref}[1]{{\color{cblue}#1}}
 and contains one tikzpicture.
% Compile with pdflatex; build_diagrams.py then crops each page to its content.
\documentclass[11pt]{article}
\usepackage[a3paper,landscape,margin=10mm]{geometry}
\usepackage[T1]{fontenc}
\usepackage{tgheros}                 % Neudata body font (as in neudata.cls)
\renewcommand{\sfdefault}{qhv}
\renewcommand{\familydefault}{\sfdefault}
\usepackage[scaled=0.95,varqu]{zi4}
\usepackage{tikz}
\usetikzlibrary{arrows.meta,positioning,fit,calc}
\pagestyle{empty}
\setlength{\parindent}{0pt}
\hyphenpenalty=10000
\exhyphenpenalty=10000

% Neudata palette (neu-data/.github BRAND.md); names kept so the diagram sources need no edits
\definecolor{cblue}{HTML}{0B376C}    % Neudata Blue: question boxes, links between diagrams
\definecolor{corange}{HTML}{B5652B}  % muted contrasting colour: warnings
\definecolor{cgreen}{HTML}{055F56}   % Neudata Teal: method boxes
\definecolor{clight}{HTML}{EAF2F5}   % Neudata Mist: question box fill
\definecolor{ctext}{HTML}{1F2D33}    % Neudata Ink: text
\definecolor{cgrey}{HTML}{8A99A3}    % Neudata Grey: further-learning routes
\definecolor{cpale}{HTML}{F4F8FA}
\definecolor{cnavy}{HTML}{04242F}    % Deep Navy: titles

\tikzset{
  every node/.style={text=ctext, font=\small},
  title/.style={font=\Large\bfseries, text=cnavy, anchor=west},
  subtitle/.style={font=\normalsize, text=ctext!72, anchor=west},
  q/.style={draw=cblue, fill=clight, rounded corners=3pt, line width=0.8pt,
            align=center, inner sep=6pt, text width=#1},
  q/.default=5cm,
  m/.style={draw=cgreen, fill=white, rounded corners=3pt, line width=1pt,
            align=left, inner sep=6pt, text width=#1},
  m/.default=5cm,
  opt/.style={draw=cblue, fill=white, rounded corners=3pt, line width=0.6pt,
              align=center, inner sep=5pt, text width=#1},
  opt/.default=3.2cm,
  further/.style={draw=cgrey, dashed, fill=white, rounded corners=3pt, line width=0.8pt,
                  align=left, inner sep=6pt, text=ctext, text width=#1},
  further/.default=5cm,
  warn/.style={draw=corange, fill=white, rounded corners=3pt, line width=1pt,
               align=left, inner sep=6pt, text width=#1},
  warn/.default=6cm,
  note/.style={align=left, text width=#1, font=\footnotesize, text=ctext},
  note/.default=5cm,
  band/.style={font=\normalsize\bfseries, text=cgreen, anchor=west, align=left},
  flow/.style={-{Stealth[length=2.6mm]}, line width=0.8pt, draw=ctext!75},
  flowg/.style={-{Stealth[length=2.6mm]}, line width=0.8pt, draw=cgrey, dashed},
  lab/.style={font=\footnotesize, fill=white, inner sep=2pt, text=ctext},
  code/.style={font=\ttfamily\footnotesize},
}
\newcommand{\rc}[1]{{\ttfamily\footnotesize #1}}
\newcommand{\dref}[1]{{\color{cblue}#1}}

\begin{document}
\begin{tikzpicture}[x=1cm,y=1cm]
\node[title] at (0,2.2) {D4. Associations and regression models};
\node[subtitle] at (0,1.5) {Regression estimates how the outcome changes with each predictor, holding the other predictors constant. Dashed boxes are beyond this course: routes for further learning.};

\node[q=11cm, anchor=north] (root) at (12.9,0.6) {\textbf{Relationship between an outcome and one or more predictors}\\\footnotesize Predictors (explanatory variables) can be continuous or categorical};
\node[q=6cm, below=5mm of root] (type) {\textbf{What type is the outcome?}};
\draw[flow] (root) -- (type);

\def\cw{5.9cm}
\foreach \i/\n/\t in {0/cont/Continuous, 1/bin/Binary, 2/cnt/Count, 3/oth/Other outcome types}{
  \node[opt=\cw, anchor=north] (\n) at (3.05+\i*6.6,-2.6) {\textbf{\t}};
  \draw[flow] (type.south) -- ++(0,-0.35) -| (\n.north);
}

% continuous
\node[m=\cw, below=5mm of cont] (cor) {\textbf{Correlation} (strength only)\\
Pearson r: linear association\\
Spearman rho: monotonic, uses ranks\\
\rc{cor.test(\textasciitilde{} sbp\_0 + age, data = hearts)}};
\node[m=\cw, below=4mm of cor] (slr) {\textbf{Simple linear regression}\\
Slope = difference in mean outcome per 1-unit difference in x, with 95\% CI\\
\rc{lm(sbp\_0 \textasciitilde{} age, data = hearts)}};
\node[m=\cw, below=4mm of slr] (mlr) {\textbf{Multiple linear regression}\\
Adjusted slopes; categorical predictors become indicator variables\\
\rc{lm(sbp\_0 \textasciitilde{} activity\_hr + age + sex,}\\\rc{\ \ \ data = hearts)}};

% binary
\node[m=\cw, below=5mm of bin] (lr) {\textbf{Logistic regression}\\
Models the log-odds. exp(coefficient) = odds ratio; adjusted OR (aOR) with several predictors\\
\rc{glm(high\_bp\_0 \textasciitilde{} smoker + age10 + sex,}\\\rc{\ \ \ \ family = binomial, data = hearts)}\\
Course data: smoking OR 3.39 crude, aOR 2.04 (95\% CI 1.17 to 3.54) adjusted for age and sex};
\node[note=\cw, below=3mm of lr] {Check: enough events per coefficient; linearity of continuous predictors on the log-odds scale. With a common outcome an OR is not a risk ratio.};

% count
\node[m=\cw, below=5mm of cnt] (pois) {\textbf{Poisson regression}\\
Models the log rate. exp(coefficient) = incidence rate ratio (IRR). Different follow-up times enter as an \textbf{offset}\\
\rc{glm(clinic\_visits \textasciitilde{} programme +}\\\rc{\ \ age10 + high\_bp\_0 +}\\\rc{\ \ offset(log(followup\_months)),}\\\rc{\ \ family = poisson, data = hearts)}};
\node[warn=\cw, below=4mm of pois] (od) {\textbf{Overdispersion:} variance larger than the mean. Pearson chi-square / residual df well above 1 (course data about 2) makes CIs too narrow. Use \rc{quasipoisson} or negative binomial (\rc{MASS::glm.nb}).};

% further learning
\node[further=\cw, below=5mm of oth] (ordl) {\textbf{Ordinal} (e.g.\ satisfaction 1 to 5)\\
Proportional odds (ordinal logistic) regression, e.g.\ \rc{MASS::polr}};
\node[further=\cw, below=4mm of ordl] (mult) {\textbf{Nominal, more than two categories}\\
Multinomial logistic regression, e.g.\ \rc{nnet::multinom}};
\node[further=\cw, below=4mm of mult] (surv) {\textbf{Time-to-event with censoring}\\
Kaplan-Meier curves, log-rank test, Cox regression (hazard ratio), \rc{survival} package};
\node[further=\cw, below=4mm of surv] (mixed) {\textbf{Repeated or clustered observations} (any outcome type)\\
Mixed-effects models or GEE: see \dref{D6}};

% bottom row: purpose and checks
\node[m=8.0cm, anchor=north west, minimum height=3.1cm] (expl) at (0.05,-13.1) {\textbf{Explaining (estimating an association)}\\
Choose adjustment variables from a causal diagram, not from p-values. Adjust for \textbf{confounders} (common causes of exposure and outcome), not for \textbf{mediators} (on the path from exposure to outcome). Association is not causation.};
\node[m=8.0cm, anchor=north west, minimum height=3.1cm] (pred) at (8.75,-13.1) {\textbf{Predicting}\\
Judge the model on data it has not seen: split into training and test sets (or cross-validate). Avoid \textbf{leakage}: no test-set or future information while building the model.};
\node[m=8.0cm, anchor=north west, minimum height=3.1cm] (chk) at (17.45,-13.1) {\textbf{Checks for linear regression}\\
Independence (from the design); linearity; constant spread; approximately normal residuals. Use \rc{plot(fit)}: residuals vs fitted, Q-Q plot. The outcome itself need not be normal.};
\end{tikzpicture}
\end{document}
